\section*{Chapitre \ref{ch:g12:ph-conductivity-titrations} --- Titrages pH-métriques et conductimétriques}

\begin{solution}{exo:g12:ph-conductivity-titrations:1}
$V_E = \qty{20.0}{mL}$~; $c = 0.100 \times 20.0 / 20.0 = \qty{0.100}{mol/L}$.
\end{solution}

\begin{solution}{exo:g12:ph-conductivity-titrations:2}
À la \omterm{def:g12:ph-conductivity-titrations:ph-metric}{demi-équivalence}, \omterm{def:g9:acids-bases-ph:ph}{pH} $= pK_a$~: $pK_a = 3.75$, proche de celui de
l'acide méthanoïque
(3.74).
\end{solution}

\begin{solution}{exo:g12:ph-conductivity-titrations:3}
La phénolphtaléine~: sa \omterm{def:g12:buffers-predominance:indicator}{zone de virage}, 8.0--10.0, est située dans le
saut, qui est basique. Le rouge de méthyle (4.2--6.3) virerait vers la
\omterm{def:g12:ph-conductivity-titrations:ph-metric}{demi-équivalence}, bien trop tôt.
\end{solution}

\begin{solution}{exo:g12:ph-conductivity-titrations:4}
La réponse de la sonde dérive avec le temps et la température~: on la
règle pour qu'elle indique la bonne valeur dans deux \omterm{def:g12:buffers-predominance:buffer}{solutions tampons}
de \omterm{def:g9:acids-bases-ph:ph}{pH} connu.
\end{solution}

\begin{solution}{exo:g12:ph-conductivity-titrations:5}
\ce{H3O+} (349.8) > \ce{OH-} (198.3) > \ce{Cl-} (76.2) > \ce{Na+} (50.3).
\end{solution}

\begin{solution}{exo:g12:ph-conductivity-titrations:6}
$c = 0.050 \times 14.6 / 20.0 = \qty{0.0365}{mol/L}$.
\end{solution}

\begin{solution}{exo:g12:ph-conductivity-titrations:7}
Pentes~: 2.0, 4.5, 17.5 et \qty{2.5}{mL^{-1}}. La plus grande se situe
entre
10.0 et \qty{10.2}{mL}~: $V_E \approx \qty{10.1}{mL}$.
\end{solution}

\begin{solution}{exo:g12:ph-conductivity-titrations:8}
Avant l'\omterm{def:g11:titration:equivalence}{équivalence}, chaque \omterm{def:g9:ions:ion}{ion} \ce{OH-} ajouté fait disparaître un \omterm{def:g9:ions:ion}{ion}
\ce{H3O+} (très conducteur) et apporte un \omterm{def:g9:ions:ion}{ion} \ce{Na+} (bien moins
conducteur)~: la \omterm{def:g12:ph-conductivity-titrations:conductivity}{conductivité} diminue. Après, les \omterm{def:g9:ions:ion}{ions} \ce{Na+} et
\ce{OH-} s'accumulent sans réagir~: elle augmente.
\end{solution}

\begin{solution}{exo:g12:ph-conductivity-titrations:9}
Au départ, peu d'\omterm{def:g9:ions:ion}{ions} (l'\omterm{def:g12:ka-and-pka:strong-acid}{acide faible} réagit à peine avec l'eau)~:
\omterm{def:g12:ph-conductivity-titrations:conductivity}{conductivité} faible. Avant l'\omterm{def:g11:titration:equivalence}{équivalence}, des \omterm{def:g9:ions:ion}{ions} éthanoate et sodium
se forment à la place de \omterm{def:g7:atoms-and-molecules:molecule}{molécules} neutres~: elle augmente, lentement.
Après l'\omterm{def:g11:titration:equivalence}{équivalence}, les \omterm{def:g9:ions:ion}{ions} \ce{Na+} et \ce{OH-} s'accumulent~: elle
augmente plus vite. Deux droites croissantes, la seconde plus pentue.
\end{solution}

\begin{solution}{exo:g12:ph-conductivity-titrations:10}
Environ 2.9, 4.76 et 8.7. L'acide éthanoïque est faible~: seuls quelques
pour cent donnent des \omterm{def:g12:ka-and-pka:oxonium}{ions oxonium}, si bien que le \omterm{def:g9:acids-bases-ph:ph}{pH} part de plus haut
que le 1.0 de l'acide chlorhydrique à une concentration voisine.
\end{solution}

\begin{solution}{exo:g12:ph-conductivity-titrations:11}
Pour que le volume de \omterm{def:g11:titration:titration}{solution titrante} ajouté modifie à peine le volume
total~: les concentrations, et donc la \omterm{def:g12:ph-conductivity-titrations:conductivity}{conductivité}, varient alors selon
des droites.
\end{solution}

\begin{solution}{exo:g12:ph-conductivity-titrations:12}
Les \omterm{def:g9:ions:ion}{ions} hydroxyde réagissent d'abord avec l'\omterm{def:g12:ka-and-pka:strong-acid}{acide fort} (\ce{H3O+}),
puis avec l'\omterm{def:g12:ka-and-pka:strong-acid}{acide faible}. Le premier saut marque la fin de l'acide
chlorhydrique (son \omterm{def:g11:titration:equivalence}{volume à l'équivalence} le mesure), le second la fin
de l'acide éthanoïque (le volume entre les deux sauts le mesure).
\end{solution}

\begin{solution}{exo:g12:ph-conductivity-titrations:13}
Au départ~: $(349.8 + 76.2) \times 0.0100 = \qty{4.26}{mS/cm}$. À
l'\omterm{def:g11:titration:equivalence}{équivalence}~: $[\ce{Na+}] = [\ce{Cl-}] = 1.00/110 = \qty{0.00909}{mol/L}$,
donc $(50.3 + 76.2) \times 0.00909 = \qty{1.15}{mS/cm}$. Les deux
valeurs sont celles de la figure.
\end{solution}

\begin{solution}{exo:g12:ph-conductivity-titrations:14}
Aucune sur $V_E$~: le saut se produit au même volume, quel que soit le
décalage des lectures. Le $pK_a$ lu à la \omterm{def:g12:ph-conductivity-titrations:ph-metric}{demi-équivalence} est trop
grand de 0.2.
\end{solution}

\begin{solution}{exo:g12:ph-conductivity-titrations:15}
Le \omterm{def:g9:acids-bases-ph:ph}{pH} ne change pas au cours de cette \omterm{def:g9:ions:precipitate}{précipitation}, mais les \omterm{def:g9:ions:ion}{ions}, si~:
avant l'\omterm{def:g11:titration:equivalence}{équivalence}, chaque \omterm{def:g9:ions:ion}{ion} \ce{Ag+} ajouté fait disparaître un \omterm{def:g9:ions:ion}{ion}
\ce{Cl-} et apporte un \omterm{def:g9:ions:ion}{ion}
\ce{NO3-} de \omterm{def:g12:ph-conductivity-titrations:conductivity}{conductivité} voisine, si bien que la \omterm{def:g12:ph-conductivity-titrations:conductivity}{conductivité} reste
presque constante (les \omterm{def:g9:ions:ion}{ions} sodium étaient présents dès le départ)~;
après l'\omterm{def:g11:titration:equivalence}{équivalence}, les \omterm{def:g9:ions:ion}{ions} \ce{Ag+} et \ce{NO3-}
s'accumulent et elle augmente. Un palier, puis une droite croissante.
\end{solution}

\begin{solution}{pb:g12:ph-conductivity-titrations:1}
\textbf{1.} \qty{6}{g} d'acide éthanoïque pour \qty{100}{g} de vinaigre.

\textbf{2.} \qty{101}{g} de vinaigre contiennent $6.0 \times 101/100 =
\qty{6.06}{g}$~; $M = \qty{60.0}{g/mol}$~: \qty{0.101}{mol} dans
\qty{0.1000}{L}, soit \qty{1.01}{mol/L}.

\textbf{3.} Non dilués, \qty{10.0}{mL} demanderaient environ
\qty{101}{mL} d'hydroxyde de sodium à \qty{0.100}{mol/L}~: plus que ne
contient une burette.

\textbf{4.} Une pipette jaugée de \qty{10.0}{mL} et une fiole jaugée de
\qty{100.0}{mL}.

\textbf{5.} \ce{CH3COOH + OH- -> CH3COO- + H2O}.

\textbf{6.} $V_E = \qty{10.1}{mL}$.

\textbf{7.} $c = 0.100 \times 10.1 / 10.0 = \qty{0.101}{mol/L}$.

\textbf{8.} $10 \times 0.101 = \qty{1.01}{mol/L}$.

\textbf{9.} \omterm{def:g9:acids-bases-ph:ph}{pH} 4.76 à \qty{5.05}{mL}~: le $pK_a$ de l'acide éthanoïque.

\textbf{10.} À l'\omterm{def:g11:titration:equivalence}{équivalence}, la solution contient de l'éthanoate de
sodium, et l'\omterm{def:g9:ions:ion}{ion} éthanoate est une \omterm{def:g12:ka-and-pka:strong-acid}{base faible}.

\textbf{11.} La phénolphtaléine, dont la \omterm{def:g12:buffers-predominance:indicator}{zone de virage}, 8.0--10.0,
contient le \omterm{def:g9:acids-bases-ph:ph}{pH} à l'\omterm{def:g11:titration:equivalence}{équivalence} (environ 8.7).

\textbf{12.} Des \omterm{def:g9:ions:ion}{ions} sodium et des \omterm{def:g9:ions:ion}{ions} éthanoate.

\textbf{13.} Chaque \omterm{def:g9:ions:ion}{ion} \ce{OH-} ajouté transforme une \omterm{def:g7:atoms-and-molecules:molecule}{molécule} presque
non ionisée en
\ce{CH3COO-}, accompagné d'un \omterm{def:g9:ions:ion}{ion} \ce{Na+}~: on ajoute lentement des
\omterm{def:g9:ions:ion}{ions} modérément conducteurs.

\textbf{14.} Au-delà de l'\omterm{def:g11:titration:equivalence}{équivalence}, les \omterm{def:g9:ions:ion}{ions} hydroxyde, très
conducteurs, s'accumulent avec les \omterm{def:g9:ions:ion}{ions} sodium.

\textbf{15.} L'\omterm{def:g12:ka-and-pka:strong-acid}{acide faible} ne cède que très peu d'\omterm{def:g9:ions:ion}{ions} à l'eau.

\textbf{16.} En $V_E = \qty{10.1}{mL}$, comme lors du \omterm{def:g12:ph-conductivity-titrations:ph-metric}{titrage
pH-métrique}.

\textbf{17.} $1.01 \times 60.0 = \qty{60.6}{g/L}$.

\textbf{18.} \qty{6.06}{g} dans \qty{100.0}{mL}.

\textbf{19.} $6.06 \times 100 / 101 = \qty{6.0}{g}$ pour \qty{100}{g}.

\textbf{20.} Le vinaigre a une acidité de \qty{6.0}{\percent}~: l'étiquette dit
vrai.
\end{solution}
